\documentclass[12pt,a4paper]{article}
\usepackage[a4paper,margin=18mm]{geometry}
\usepackage[T2A]{fontenc}
\usepackage[utf8]{inputenc}
\usepackage[russian]{babel}
\usepackage{amsmath,amssymb,mathtools}
\usepackage{array,booktabs,longtable}
\usepackage{xcolor}
\usepackage{hyperref}
\usepackage{tikz}

\hypersetup{
  colorlinks=true,
  linkcolor=blue!55!black,
  urlcolor=blue!55!black
}
\setlength{\parindent}{0pt}
\setlength{\parskip}{0.55em}
\renewcommand{\arraystretch}{1.25}

% Краткое сообщение Софии: София, в этой работе стоит отдельно вернуться к задаче №6: решение там не начато. Ключевой ход — заметить подобие двух прямоугольных треугольников по краям квадрата; из него получается связь между площадью квадрата, AM и BN. Ниже этот переход разобран пошагово.

\begin{document}

\begin{titlepage}
\thispagestyle{empty}
\begin{center}
\vspace*{0.24\textheight}
{\Huge\bfseries Ревью домашней работы\par}
\vspace{2.2cm}
{\Large Автор: Лёвин Артём Александрович\par}
\vspace{1.0cm}
{\large 10 сентября 2026 года\par}
\end{center}
\vfill
\begin{center}
\begin{tikzpicture}[x=1cm,y=1cm]
\draw[blue!55!black, line width=0.7pt]
  (0,0) .. controls (2.0,0.55) and (4.0,-0.55) .. (6.0,0)
        .. controls (8.0,0.55) and (10.0,-0.55) .. (12.0,0);
\end{tikzpicture}
\end{center}
\end{titlepage}

\newpage
\section*{Сводная таблица}
\footnotesize
\setlength{\tabcolsep}{4pt}
\begin{longtable}{>{\centering\arraybackslash}p{0.08\linewidth}p{0.18\linewidth}p{0.26\linewidth}p{0.14\linewidth}p{0.15\linewidth}}
\toprule
\textbf{№ задания} & \textbf{Тема} & \textbf{Суть ошибки} & \textbf{Ваш ответ} & \textbf{Правильный ответ} \\
\midrule
\endfirsthead
\toprule
\textbf{№ задания} & \textbf{Тема} & \textbf{Суть ошибки} & \textbf{Ваш ответ} & \textbf{Правильный ответ} \\
\midrule
\endhead
\hyperref[task:6]{6} & Подобие прямоугольных треугольников; квадрат в прямоугольном треугольнике & Решение не выполнено: не использована связь, которая следует из подобия треугольников по краям квадрата. & не указан & $BN=9$ \\
\bottomrule
\end{longtable}
\normalsize

\newpage
\null\vfill
\begin{center}
{\LARGE\bfseries Разбор ошибок}
\end{center}
\vfill\null
\newpage

\phantomsection\label{task:6}
\section*{Задание 6}

\subsection*{Текст задания}
В прямоугольном треугольнике $ABC$ квадрат $KLMN$ расположен так же, как в задаче 1: $M,N\in AB$, $L\in AC$, $K\in BC$. Площадь квадрата равна $54$, а $AM=6$. Найдите $BN$.

\subsection*{Ваше решение}
Записано: «не поняла». Математическое решение дальше не начато.

\subsection*{Ваш ответ}
Не указан.

\subsection*{Ошибка}
\textcolor{red}{Ошибка:} решение не выполнено; пропущено ключевое применение подобия двух прямоугольных треугольников, прилегающих к квадрату слева и справа.

\subsection*{Где именно возникает ошибка}
Последнего корректного шага решения здесь нет, потому что после условия не записано ни одного математического перехода.

Первый необходимый, но пропущенный шаг --- обозначить сторону квадрата через $x$ и связать два крайних прямоугольных треугольника $ALM$ и $KNB$ по углам. Именно это даёт уравнение, в котором одновременно участвуют $AM$, $BN$ и сторона квадрата.

\subsection*{Почему это неверно}
Само отсутствие вычислений не означает, что данных недостаточно. Конфигурация полностью определяет нужную связь.

Поскольку $KLMN$ --- квадрат, его стороны $LM$ и $KN$ перпендикулярны $AB$ и равны между собой. Поэтому треугольники $ALM$ и $KNB$ прямоугольные.

Пусть $\angle A=\alpha$. Так как $\triangle ABC$ прямоугольный при $C$, то
\[
\angle B=90^\circ-\alpha.
\]
В треугольнике $ALM$ угол при $M$ равен $90^\circ$, значит
\[
\angle ALM=90^\circ-\alpha.
\]
В треугольнике $KNB$ угол при $N$ также равен $90^\circ$, а угол при $B$ равен $90^\circ-\alpha$, поэтому
\[
\angle BKN=\alpha.
\]
Следовательно,
\[
\triangle ALM\sim\triangle KNB
\]
по двум углам.

При этом соответствуют стороны
\[
AM\leftrightarrow KN,\qquad LM\leftrightarrow BN.
\]
Поэтому
\[
\frac{LM}{BN}=\frac{AM}{KN}.
\]
Так как $LM=KN=x$, получаем
\[
\frac{x}{BN}=\frac{AM}{x},
\]
то есть
\[
x^2=AM\cdot BN.
\]
Но $x^2$ --- это площадь квадрата. Значит, для данной конфигурации
\[
S_{\text{кв}}=AM\cdot BN.
\]
Именно эта связь позволяет найти $BN$ напрямую.

\subsection*{Ключевая идея}
\textcolor{blue!60!black}{Ключевая идея:} не нужно находить стороны большого треугольника $ABC$. Достаточно увидеть подобие двух маленьких прямоугольных треугольников по краям квадрата. Из подобия сразу следует
\[
S_{\text{кв}}=AM\cdot BN.
\]

\subsection*{Исправление от последнего правильного шага}
Так как записанного решения нет, исправление начинается с первого содержательного шага.

Обозначим сторону квадрата через $x$. Тогда
\[
x^2=S_{\text{кв}}=54.
\]
Из подобия $\triangle ALM\sim\triangle KNB$:
\[
\frac{LM}{BN}=\frac{AM}{KN}.
\]
Подставляем $LM=KN=x$ и $AM=6$:
\[
\frac{x}{BN}=\frac{6}{x}.
\]
Перемножаем крест-накрест:
\[
x^2=6BN.
\]
Но $x^2=54$, поэтому
\[
54=6BN.
\]
Отсюда
\[
BN=9.
\]

\subsection*{Полное правильное решение}
Пусть сторона квадрата $KLMN$ равна $x$.

Из площади квадрата:
\[
x^2=54.
\]
При необходимости можно найти саму сторону:
\[
x=\sqrt{54}=3\sqrt6,
\]
причём берём только положительное значение, так как длина стороны положительна.

Пусть $\angle A=\alpha$. Так как $\angle C=90^\circ$, то
\[
\angle B=90^\circ-\alpha.
\]
Поскольку $LM\perp AB$ и $KN\perp AB$, треугольники $ALM$ и $KNB$ прямоугольные. В них
\[
\angle A=\angle BKN=\alpha,
\]
а также
\[
\angle AML=\angle KNB=90^\circ.
\]
Следовательно,
\[
\triangle ALM\sim\triangle KNB.
\]
Из подобия:
\[
\frac{LM}{BN}=\frac{AM}{KN}.
\]
Так как $LM=KN=x$ и $AM=6$, имеем
\[
\frac{x}{BN}=\frac{6}{x}.
\]
Отсюда
\[
x^2=6BN.
\]
Подставляем $x^2=54$:
\[
54=6BN,
\]
\[
BN=\frac{54}{6}=9.
\]

\subsection*{Проверка}
Полученное значение положительно, что соответствует длине отрезка. Проверим основную связь:
\[
AM\cdot BN=6\cdot 9=54=S_{\text{кв}}.
\]
Равенство выполняется, значит найденное значение согласуется с подобием треугольников и площадью квадрата.

\subsection*{Итог}
\textcolor{teal!60!black}{Итог:} $BN=9$. Для задач такой конфигурации полезно запомнить не как готовую формулу, а как следствие подобия: квадрат создаёт два прямоугольных треугольника, из которых получается связь $S_{\text{кв}}=AM\cdot BN$.

\subsection*{Задача на закрепление}
В прямоугольном треугольнике квадрат расположен так же, как в задаче 1. Площадь квадрата равна $80$, а $AM=5$. Найдите $BN$.

\vfill
\begin{center}
\small Лёвин Артём Александрович эксклюзивно для Софии Кальней
\end{center}

\end{document}
